% !TeX root = ../main.tex
\chapter{So sánh kết quả (Comparison of Results)}
\label{chap:comparison}

\section{Bảng tổng hợp}
Dưới đây là bảng tổng hợp kết quả của các mô hình (cơ sở so sánh từ bảng \texttt{results/model\_comparison.csv}):

\begin{table}[H]
\centering
\caption{So sánh Baseline và các phương pháp cải thiện}
\label{tab:model-comparison}
\resizebox{\textwidth}{!}{%
\begin{tabular}{llcccc}
\toprule
\textbf{Model} & \textbf{Key Parameter} & \textbf{Train Accuracy} & \textbf{Test Accuracy} & \textbf{Tree Depth} & \textbf{Leaves} \\
\midrule
Baseline         & default            & 100.0\%        & 91.23\%       & 7          & 19     \\
Max Depth        & max\_depth=4       & 98.68\%        & 93.86\%       & 4          & 11     \\
Min Samples Leaf & min\_samples\_leaf=3 & 98.02\%        & 91.23\%       & 6          & 13     \\
Pruning          & ccp\_alpha=0.005934 & 97.36\%        & 93.86\%       & 3          & 6      \\
\bottomrule
\end{tabular}}
\end{table}

\section{Xác định Winner (Mô hình tốt nhất)}
Mô hình tốt nhất là \textbf{Max Depth (\texttt{max\_depth=4})}. Dù Max Depth và Pruning chia sẻ chung vị trí dẫn đầu về Test Accuracy (93.86\%), Max Depth được chọn do có chỉ số \textbf{Malignant Recall (độ nhạy với khối u ác tính)} là 92.86\%, cao hơn mức 90.48\% của Pruning. Trong lĩnh vực y khoa, việc bỏ sót chẩn đoán một bệnh nhân ung thư nguy hiểm hơn rất nhiều so với chẩn đoán nhầm (đối với lành tính), vì vậy ưu tiên Recall cho lớp Malignant là hợp lý.

\section{Đánh giá Trade-off}
Baseline đánh đổi khả năng dự đoán thực tế để đạt 100\% trên dữ liệu đã biết (độ phức tạp cao). Trong khi đó, Max Depth và Pruning chấp nhận hy sinh vài phần trăm độ chính xác của tập huấn luyện (Train Acc tụt xuống dưới 100\%) để làm giảm mức độ phức tạp. Sự đánh đổi này mang lại lợi ích khổng lồ, khi mô hình thực sự đã cải thiện được khả năng ``hiểu'' dữ liệu thay vì ``học thuộc lòng''.
